\section{Introduction}
\done\todo{
Introduce the problem: A bottleneck to improve automated vulnerability identification and fix in software is the lack of labeled data~\cite{dil2025,chen2024improvingdatacurationsoftware}.}

With technology constantly growing, so is the need to maintain its software against many possible vulnerabilities. Manually scanning code to detect weaknesses has become increasingly impractical, as a result, algorithms have been developed to detect vulnerabilities; one implementation is through the use of Machine Learning (ML) and Artificial Intelligence (AI) models, used to automatically scan and detect when code has a vulnerability \cite{mirsky2023vulchecker, fu2022linevul, chen2024improvingdatacurationsoftware} \cite{dil2025}. These AI/ML models must be trained with corpora of program code instances that are already known to have and to not have vulnerability bugs, so that the model can learn to recognize the associated patterns.

A major issue with this process is the collection of reliable training data. Databases like the National Vulnerability Database (NVD) and Commmon Vulnerabilities Exposures (CVE) list, both publicly available government-run/sponsored databases, can be used to obtain code with vulnerabilities. However, they have been proven to provide poor capabilities for training AI/ML, as much of the data is incomplete or low quality which makes it less likely to use for a practical prediction model \cite{zhang2011empirical}. The primary reason is that these databases are not designed toward training AI/ML models, rather for quickly reporting newly discovered vulnerabilities. Because of this focus on speed, much of the patch data can be lacking in necessary details required for training these models. This problem can be circumvented through filtering: evaluating each data point individually and removing those with issues. However, this task is too labor intensive and leads to a small data set. To scale the amount of data while also maintaining quality, automation of this collection is needed. This requires training models to filter through and to identify and collect the most useful data. Without any filtering, any data not evaluated causes noise and can lead to false negatives or false positives in vulnerability detection. Several automatic curation methods have been previously proposed. One implementation uses uncertainty quantification, a technique that measures outcomes despite unknown variables, to find data sets that have low aleatoric but high epistemic uncertainty~\cite{chen2024improvingdatacurationsoftware}. The resulting training data is high in both quality and usefulness. Another implementation uses LLMs with certain prompting strategies to best leverage the model in detecting vulnerabilities~\cite{dil2025}.

An issue with these previous works is that these algorithms are evaluated solely on existing, precurated datasets. This does not test the model's ability of obtaining new code instances, which in turn does not test the model's applicability in detecting code in the real world. This is a crucial step, as a model that is trained on unrealistic labeled data will have decreased performance ~\cite{jimenez2019importance}, due to the model not being equipped to handle real world code.

To improve upon this, we propose to a more practical method of collecting quality data. Since using human labeled data in conjunction with auto labeled data has been shown to improve performance by up to 22\% \cite{le2024automatic} in vulnerability detection, this method leverages both the UQ model and the LLM to create a small enough subset of quality data feasible for human annotation. Human annotation results in the most valuable training data, but the expertise it requires is not widely available and is also costly. Therefore, the goal is to maximize the ratio between the benefit and cost of human annotation by decreasing the amount of annotation needed. A previous study has shown that using as few as 125 human generated data points, with the rest being synthetic data, has improved performance of language models \cite{ashok2024little}. Considering this study, we want to analyze whether there is also significant improvement to performance of the vulnerability detection model if we add a set of around 100 human annotated data points. If this method is able demonstrate that we can just label a small number of functions, and have it greatly improve the performance, we can greatly reduce costs and streamline the process of securing software, and further mitigate the risk of security breaches due to software vulnerabilities.
%software vulnerability labeling, vulnerability database


\done\todo{
Discuss motivation of the problem}
\done\todo{Discuss an overview of the research approach, the cost-effective way to identify labeled data.}
